\section{Introduction}
\label{sec:intro}

Evaluating a spherical-harmonic gravity field to degree $N$ costs
$\mathcal{O}(N^2)$, and an adaptive-step integrator makes tens of thousands of
force calls per propagated day. On an eccentric orbit the required degree is
not constant along the arc, so the choice of $N$ is not one number but a
schedule, and choosing that schedule under a declared compute budget is a
resource-allocation problem rather than an accuracy question.

The natural way to allocate is from the local truncation defect: spend degrees
where the omitted acceleration is largest. That is what an attenuation
threshold, a Kaula-weighted tail criterion, or a radius-driven rule all do in
different guises, and it is what a Lagrangian treatment of a budgeted
separable allocation formalizes \citep{everett1963generalized,
shoham1988efficient}. It also produces the wrong ranking. Held to one gravity
budget on two independent coverage designs, a schedule concentrated near
perilune lowers the trajectory-averaged truncation force defect by a median
factor of about five and still gives the larger seven-day position error
\citep[\ph{self-citation of the previous paper}]{atallah2022radial}. The
mechanism is not a sampling artefact: a first-order displacement is a
\emph{signed} integral of the defect against the state-transition matrix, and a
defect norm is a sum of magnitudes, so the two agree only when the defect is
incoherent, which on these arcs it is not.

\dnote{The self-citation above is the previous paper. Once it has a DOI,
replace the placeholder and re-check that every claim attributed to it here is
one that paper actually makes at the confidence level it makes it.}

That leaves an obvious question unanswered, and this paper answers it. If the
force-level criterion is the wrong objective, what does the right one give, how
much is it worth, and can any of it be obtained without the reference-field
access the benchmark is granted? Three questions organize the paper.

\begin{description}
  \item[RQ1] Under trajectory-aware information, how much accuracy is
    \emph{attainable} at a fixed budget, relative to the constant degree and to
    the best radius-driven schedules?
  \item[RQ2] How much of that comes from weighting each epoch by its
    sensitivity, and how much from the signed cancellation that a magnitude
    criterion discards?
  \item[RQ3] How much survives without the reference trajectory and without
    pointwise reference-field evaluations along it --- the controller keeping
    only global, offline model metadata such as the degree variances?
\end{description}

Three things make the answers non-obvious, and each is a contribution.

\paragraph{The objective is non-separable.} Writing the trajectory error in the
form the dynamics impose puts the norm \emph{after} the signed integral, which
couples every pair of epochs through cancellation. The standard multiplier
treatment therefore does not apply, and neither does the obvious repair of
reweighting the per-epoch defect by its sensitivity: that corrects the weight
while still summing magnitudes. Section~\ref{sec:problem} states the problem
and Section~\ref{sec:algorithm} shows that it is nonetheless cheap to solve
well, because the discretized coupling matrix depends on the epoch pair only
through $\max(i,k)$. A full coordinate sweep costs
$\mathcal{O}(M\lvert\mathcal{N}\rvert)$, and a Frank--Wolfe bound on the
convex-hull relaxation \citep{frank1956algorithm,jaggi2013revisiting}
certifies how far the result sits from the relaxed optimum. Nothing is
propagated to obtain it. The degree is a decision variable on a coarse decision
grid rather than on the fine grid the objective is accumulated over, which
costs nothing in the solve and matters for the comparison: a benchmark allowed
to switch faster than any deployable controller would earn part of its margin
from resolution rather than from information.

\paragraph{Deployment needs a direction, not a magnitude.} A compact tail rule
predicts how much acceleration a truncation omits; it says nothing about which
way it points. Cancellation is exactly the part that needs the direction, so a
controller built on a tail magnitude alone can reach the
sensitivity-weighted level and no further. Section~\ref{sec:controller} shows
that tabulating the direction is not a cheap way out either: the omitted field
varies on the model's own resolution scale $\pi r/N$, which at degree 300 is
under twenty kilometres of ground track, so a surrogate coarser than that scale
cannot represent it and one fine enough to do so carries a comparable number of
degrees of freedom. What is affordable is to \emph{probe} it, by evaluating the
leading omitted bands --- one shared stack serving every candidate degree ---
and to complete the amplitude from the model's measured degree-variance
spectrum.

\paragraph{The plan and the probe have different lifetimes.} The cancellation
target is an accumulated quantity and is stable along an arc; the local omitted
vector is field texture and is not. The same resolution scale that rules out a
tabulated direction also says how far ahead one may look: over a single
decision interval the vehicle's own state prediction is far more accurate than
$\pi r/N$, so probing forward along it is admissible, while over a seven-day
pilot arc it is not. A controller that plans offline on a cheap pilot arc and
probes forward one interval at a time respects that split, reads neither the
reference arc nor the field along it, and is the object the
campaign of Section~\ref{sec:campaign} evaluates.

Throughout, the budget is a common \emph{ceiling} on realized total quadratic
work, with the controller's own overhead charged against it, and not a nominal
per-call count. A ceiling rather than an equality because the objective is not
monotone in degree (Section~\ref{sec:algorithm}), so realized work is reported
beside every error rather than assumed identical. The
previous paper measured that the two differ by a median \ph{pct} and by a
factor of \ph{value} in the compute-starved regime, which is precisely the
regime where a budget matters, so the weaker definition is not used here to
carry a verdict.

\dnote{If the campaign lands in the amber band --- benchmark gap real,
controller cannot capture it --- this introduction inverts: the three
contributions become (i) the non-separable formulation and its certificate,
(ii) the resolution-scale argument as a measured \emph{cost} statement about
surrogates for the omitted direction, and (iii) a measured account of what deployable rules leave
on the table. The formulation and the decorrelation argument survive either
way; only the third paragraph changes.}
